\begin{abstract}
Embodied question answering (EQA) requires a robot to explore an unknown environment, build a semantic model of what it sees, and answer natural-language questions with sufficient confidence to know when to keep looking.
We present \textbf{Dynagraph}, which keeps \textbf{DynaMem's voxel update rules} and adds \textbf{GraphEQA graph-creation steps} on the same perception stream.
Each voxel records when it was last seen and a pointer to the supporting observation (pose, robot configuration, images); object-centric graphs are \textbf{grounded in this map}, so the system can detect when an object is \textbf{no longer observed} and retire stale hypotheses.
Dynagraph adds \textbf{spatial merge} and \textbf{staleness pruning} so memory stays compact and current as the robot moves and the scene changes.
We implement Dynagraph in the open-source \texttt{emet} stack (\texttt{emet run dynagraph}) and evaluate on \textbf{Habitat} (HM-EQA, via a dedicated harness) and on \textbf{MuJoCo simulation} (Robocasa kitchens, MolmoSpaces iTHOR, multiple robot embodiments).
Our primary evaluation compares two controllers on the \emph{same} Dynagraph memory: a \textbf{classic} GraphEQA planning loop and an \textbf{agentic} tool controller that retrieves place evidence cards, navigates / explores, and uses cheap detector proposals only as high-recall candidates while the answer is VLM-assessed before submission.
On an anti-overfit holdout of eight HM-EQA questions (matched same-policy, router off), agentic answers \textbf{8/8} vs.\ classic \textbf{5/8} with roughly $3\times$ fewer planning steps ($18.3$ vs.\ $60.3$).
On a letter-balanced 32-question slice, agentic reaches \textbf{16/32} vs.\ classic \textbf{9/32} with significantly fewer mean planning steps (cross-policy composite; Wilcoxon $p{\approx}2{\times}10^{-5}$ on steps, McNemar $p{\approx}0.09$ on letters).
A full 113-question comparison against GraphEQA's API-VLM baselines with the current navigation stack is in progress; preliminary local-VLM slices trail those baselines and are at least competitive with our static-graph baseline on matched slices (+3\,pp on a letter-balanced 32 at 3\,B), and the agentic loop also lowers search cost.
Other benchmark tables (OVMM find-phase, SQA3D, dynamic exploration) are pending large-scale sim runs.
\end{abstract}
